\par\SourceLeaf{493}%
In the meantime I have tried in several ways to obtain information that will be important for the Hanoverian degree survey in several respects. The survey takes on greater importance because it continues the extensive Danish survey on the one hand and is capable of a much more extensive continuation southward on the other. I am therefore glad to be able to say now that a substantial part of the latter operations, which previously could be described only as possible and desirable, has already in fact been carried out. The Royal Prussian Government intends to have several principal provinces of the Prussian monarchy surveyed and, in accordance with principles now recognized as the only admissible ones, is first laying the foundation for this survey with a large triangulation. It has already had the triangulation \textit{actually} carried out that connects the Prussian Rhine provinces to the older ones and extends through Nassau, Electoral Hesse, Thuringia, and so on, as far as Halle. The Royal Prussian Lieutenant General von \textsc{Müffling} communicated the whole triangulation to me, and after careful examination of my own I was convinced that it had been carried out with such care and accuracy that it can rank alongside the triangulations in the best degree surveys and even surpasses some of them. Some of these triangles pass very close to the southern borders of the Kingdom of Hanover. If, therefore, my own future measurements are duly connected to them, for which Lieutenant General von \textsc{Müffling} has already provided the essential information, the Danish-Hanoverian degree survey will thereby already be extended to the Seeberg observatory near Gotha and to the Inselsberg. To the southwest, a connection with the observatories of Mannheim and Paris will also be achieved by this means and through the Darmstadt and French surveys.

I have also obtained, in several ways and partly direct from the Dépot de la guerre, several important though incomplete pieces of information on the measurements carried out in the Electorate of Hanover during the French occupation in the years 1803—1805 by the French colonel \textsc{Epailly}. Earlier efforts to obtain full particulars by diplomatic means proved fruitless, and I still have some hope of being sent \textit{all} the triangles. If a general triangulation of the entire Kingdom of Hanover is ever decided upon, for which the degree survey from Hamburg to the southern border of the kingdom will provide the soundest basis, access to those measurements of \textsc{Epailly} could undoubtedly prove exceedingly useful. What has been communicated to me at least enables me, in arranging my own triangles, to take account of such a possible use in advance and to prepare for it to some extent. — Among the documents concerning \textsc{Epailly}’s survey that have come into my hands there is also a report by Colonel \textsc{Epailly} himself, addressed to General \textsc{Samson}. This has indeed made me aware in advance of the great \textit{difficulties} of the triangulations, arising from various local circumstances, particularly in the Lüneburg region. The French geodesist himself found those difficulties almost insurmountable even in the ‘pays conquis’, and I can hope to overcome them only with strong support from the government and the authorities concerned.

I have also made suitable arrangements over the course of this winter for various other requirements of the coming surveys. The large \textsc{Reichenbach} theodolite that is expected is intended for the final measurement of the angles with the utmost precision. But finding the stations, reconnoitering the points visible from each of them, and making the preliminary measurements require another theodolite that is easy to transport, quick to set up and easy to handle, quite apart from the fact that the principal instrument must be treated with the greatest care. For that purpose
